%=====================================================================
\section{Conclusion}
\label{sec:conclusion}
The completed citation-network experiments show that the importance of time
depends on the question. A final aggregate reveals degree heterogeneity and
connectivity, while preferential attachment, ageing and the growth law require
historical degrees, publication ages or cumulative snapshots. In the selected
AMiner database corpus, sublinear degree preference and age-dependent citation
choice coexist with superlinear edge growth, yet the sampled effective diameter
does not shrink. WSB compactly describes selected citation histories, but its
mean-absolute-error advantage over a simple linear forecast is not established.
Thus descriptive fit, temporal mechanism and predictive value should be assessed
separately. PS/CM experiments remain pending, so conclusions about their dynamics
and a cross-network ranking of temporal sensitivity await measured evidence.
Future work includes broader corpus coverage, more datasets per type,
finer-resolution citation dates and predictive tests under alternative censoring
and missing-reference assumptions.
